% Unique practice relocated from topic 2.2. Original board papers not claimed.
\begingroup
\chTwoReadableTopicStyle
\subsection{ট্রান্সমিশন: স্থানান্তরিত ও প্রয়োগমূলক অনুশীলন}
\label{ch2:transmission-practice}
এগুলো নিজস্ব ধারণাভিত্তিক অনুশীলন। Walkie-talkie, ফোন ও একজন প্রাপকের MCQ আগের তালিকায় আছে; এখানে আবার দেওয়া হয়নি। Frame Efficiency ও Bit Synchronization-এর অতিরিক্ত MCQ-ও আগের ধারণাভিত্তিক তালিকায় আছে।
\begin{multicols}{2}\begin{enumerate}
\item\chTwoMCQBlock{প্যারালাল ট্রান্সমিশনে ডেটা কীভাবে যায়?
\par ক. বিট পরপর যায় \quad খ. একই সময়ে একাধিক বিট যায় \par গ. শুধু তারবিহীন পথে যায় \quad ঘ. শুধু অডিও পাঠায় }
\item\chTwoMCQBlock{সিরিয়াল ট্রান্সমিশনে ডেটা কীভাবে যায়?
\par ক. সব বিট একসঙ্গে যায় \quad খ. বিট পরপর যায় \par গ. শুধু ব্রডকাস্ট করে \quad ঘ. মাধ্যম ছাড়াই যায় }
\item\chTwoMCQBlock{প্রচলিত অ্যাসিনক্রোনাস ক্যারেক্টার ফ্রেমের সীমানা চেনার উপাদান কোনটি?
\par ক. স্টার্ট ও স্টপ বিট \quad খ. শুধু প্যারিটি \par গ. IP বিট \quad ঘ. ক্লাউড বিট }
\item\chTwoMCQBlock{সরাসরি অডিও ও ভিডিওতে নিয়মিত সময়ের ডেটাপ্রবাহ কোন ধারণার সঙ্গে সম্পর্কিত?
\par ক. Asynchronous \quad খ. Isochronous \par গ. Simplex \quad ঘ. শুধু Broadcast }
\item\chTwoMCQBlock{8 ডেটা, 1 স্টার্ট, 1 প্যারিটি ও 2 স্টপ বিট থাকলে ফ্রেমে মোট কত বিট হয়?
\par ক. 10 \quad খ. 11 \quad গ. 12 \quad ঘ. 13 }
\item\chTwoMCQBlock{100টি 8-বিট ক্যারেক্টারের ব্লকে মোট বাড়তি নিয়ন্ত্রণতথ্য 40 বিট। কার্যদক্ষতা প্রায় কত?
\par ক. 80\% \quad খ. 72.73\% \par গ. 100\% \quad ঘ. 95.24\% }
\item\chTwoMCQBlock{মাল্টিকাস্ট গ্রুপে পাঠানো থেকে কোন সিদ্ধান্ত নিশ্চিত করা যায় না?
\par ক. নির্বাচিত প্রাপকদের দল রয়েছে \par খ. প্রাপক একজনের বেশি হতে পারে \par গ. একই সময়ে দুই দিকে ডেটা পাঠানো যায় \par ঘ. সবার কাছে পাঠানো থেকে এটি আলাদা শ্রেণি }
\item\chTwoMCQBlock{প্রেরকের প্রধান স্মৃতিতে কয়েকটি ক্যারেক্টার জমিয়ে ব্লক তৈরি করা হচ্ছে। এই অধ্যায়ের প্রচলিত ক্যারেক্টার ও ব্লকভিত্তিক শ্রেণিতে কোন সময়সমন্বয় পদ্ধতি নির্দেশিত হয়?
\par ক. Asynchronous \quad খ. Synchronous \par গ. Simplex \quad ঘ. Broadcast }
\end{enumerate}\end{multicols}
\textbf{উত্তরমালা:} ১. খ, ২. খ, ৩. ক, ৪. খ, ৫. গ, ৬. ঘ, ৭. গ, ৮. খ।

\par\Needspace{7\baselineskip}
\textbf{CQ ১: কলেজের অনলাইন ক্লাস}\par
শিক্ষক সরাসরি ভিডিও ক্লাস নেন। শিক্ষার্থীরা চ্যাটে প্রশ্ন করে এবং ক্লাসের রেকর্ডিং সার্ভারে আপলোড হয়। ব্যবস্থায় শিক্ষক ও শিক্ষার্থীর একই সময়ে কথা বলা ও শোনার সুযোগ রয়েছে।
\begin{enumerate}
\item ক. সিরিয়াল ট্রান্সমিশন কী?
\item খ. অ্যাসিনক্রোনাস ট্রান্সমিশনে স্টার্ট ও স্টপ বিট প্রয়োজন কেন?
\item গ. সরাসরি ক্লাসে ডেটা পাঠানোর কোন সময়সংক্রান্ত ধারণা বেশি প্রয়োজন? ব্যাখ্যা করো।
\item ঘ. সরাসরি পারস্পরিক যোগাযোগের সঙ্গে ফুল-ডুপ্লেক্সের সম্পর্ক বিশ্লেষণ করো।
\end{enumerate}

\par\Needspace{7\baselineskip}
\textbf{CQ ২: প্রতিষ্ঠানের সভা ও প্রশিক্ষণ}\par
অনলাইন সভায় একই সময়ে কথা বলা ও শোনা যায়। ফাইল সার্ভার থেকে একজন নির্দিষ্ট ব্যবহারকারী ডাউনলোড করেন। মাল্টিকাস্ট-সমর্থিত সরাসরি প্রশিক্ষণ একটি Group Address-এ নিবন্ধিত প্রাপকদের কাছে পাঠানো হয়।
\begin{enumerate}
\item ক. ডেটা প্রবাহের দিকভিত্তিক ট্রান্সমিশন মোড কী?
\item খ. দীর্ঘ দূরত্বে সিরিয়াল ট্রান্সমিশন সুবিধাজনক কেন?
\item গ. উদ্দীপকে ফুল-ডুপ্লেক্স ও ইউনিকাস্ট শনাক্ত করার বৈশিষ্ট্যগুলো ব্যাখ্যা করো।
\item ঘ. প্রশিক্ষণে মাল্টিকাস্ট এবং নিয়মিত সময়ের ডেটাপ্রবাহের প্রয়োজন আলাদাভাবে বিশ্লেষণ করো।
\end{enumerate}
\endgroup
